\documentclass[12pt]{article}
\usepackage[T1]{fontenc}
\usepackage[utf8]{inputenc}
\usepackage{lmodern}
\usepackage{textcomp}
\usepackage{enumitem}
\usepackage{geometry}
\geometry{margin=1in}
\begin{document}
\begin{center}
Answer Key - Economics - Undergraduate Level Assessment 13 \
\small (Answers are ordered exactly as in the question paper.)
\end{center}

\section*{Multiple Choice}
\begin{enumerate}[label=\arabic*.]
\item \textbf{Answer:} A
\\ \textit{Options:} A) Product differentiation and loss of perfect substitutability; B) Barriers to entry and exit in the market; C) High consumer loyalty to a single brand; D) Lack of government regulation; E) Constant production costs regardless of output levels
\\ \textit{Note:} The correct answer highlights that product differentiation in monopolistically competitive industries leads to variations in products, which results in imperfect substitutability and complicates the definition of these industries compared to perfectly competitive markets.
\item \textbf{Answer:} D
\\ \textit{Options:} A) Automobiles depreciate faster; B) Stocks and bonds are physically smaller; C) Automobiles are more expensive; D) Highly efficient and fast markets exist for stocks and bonds
\\ \textit{Note:} The correct answer highlights that the existence of highly efficient and fast markets for stocks and bonds allows for quick buying and selling, making them more liquid than an automobile, which typically requires more time and effort to sell.
\item \textbf{Answer:} A
\\ \textit{Options:} A) a decrease in the price of coffee.; B) an increase in consumer surplus due to lower coffee prices.; C) a persistent shortage of coffee in the market.; D) a decrease in the demand for coffee.; E) a decrease in consumer surplus due to higher coffee prices.
\\ \textit{Note:} The correct answer is a decrease in the price of coffee because a price floor, which sets a minimum price above the equilibrium, leads to a surplus of coffee, causing downward pressure on prices as suppliers attempt to sell excess inventory.
\item \textbf{Answer:} D
\\ \textit{Options:} A) an elimination of consumer surplus.; B) a decrease in the demand for the good.; C) a decrease in producer surplus.; D) a persistent surplus of the good.; E) an increase in total welfare.
\\ \textit{Note:} A price floor, set above the equilibrium price, leads to a situation where the quantity supplied exceeds the quantity demanded, resulting in a persistent surplus of the good.
\item \textbf{Answer:} D
\\ \textit{Options:} A) the demand for goods and services decreases.; B) the stock market crashes.; C) bond prices fall.; D) the nominal interest rate falls.; E) the unemployment rate rises.
\\ \textit{Note:} Households demand more money as an asset when the nominal interest rate falls because lower interest rates reduce the opportunity cost of holding cash, making it more attractive to maintain liquid assets instead of investing in interest-bearing assets.
\end{enumerate}

\section*{True / False}
\begin{enumerate}[label=\arabic*.]
\setcounter{enumi}{5}
\item \textbf{Answer:} False
\\ \textit{Options:} A) True; B) False
\\ \textit{Note:} The statement is false because, in an ARMA(p,q) model, while the autocorrelation function (ACF) becomes zero after q lags, the partial autocorrelation function (PACF) generally does not become zero after p lags, reflecting the influence of the autoregressive component.
\item \textbf{Answer:} True
\\ \textit{Options:} A) True; B) False
\\ \textit{Note:} The kinked demand curve illustrates that at the kink point, the demand is elastic for price increases (as consumers will switch to substitutes) and inelastic for price decreases (since competitors are likely to match price cuts), leading to a discontinuous demand function.
\item \textbf{Answer:} True
\\ \textit{Options:} A) True; B) False
\\ \textit{Note:} When price discrimination occurs, firms charge different prices to different consumers based on their willingness to pay, resulting in the firm capturing a larger portion of the consumer surplus, which consequently reduces the overall consumer surplus in the market.
\item \textbf{Answer:} True
\\ \textit{Options:} A) True; B) False
\\ \textit{Note:} Investment demand most likely increases during a stock market downturn because investors tend to shift their funds from riskier assets to safer investments, such as bonds or cash equivalents, to preserve capital and avoid potential losses.
\item \textbf{Answer:} False
\\ \textit{Options:} A) True; B) False
\\ \textit{Note:} Required reserves can be held in various forms, including as deposits at the Federal Reserve, not just in a bank's vault.
\end{enumerate}

\section*{Long Form Response}
\begin{enumerate}[label=\arabic*.]
\setcounter{enumi}{10}
\item \textbf{Answer:} The late nineteenth-century single-tax movement aimed to reform the taxation system by eliminating all taxes except for a tax on land. This movement was inspired by the Ricardian theory of land rents, which posits that land value is derived from its location and societal demand rather than individual effort. The economic rationale behind this approach is that taxing land would capture the unearned income generated by landowners, promoting equitable wealth distribution and discouraging speculation. By shifting the tax burden away from labor and capital, proponents believed it would incentivize productive use of resources and stimulate economic growth. Ultimately, the single-tax movement sought to address social inequalities while promoting a more efficient allocation of land and resources in the economy.
\\ \textit{Note:} The answer correctly identifies the single-tax movement's goal of reforming taxation by focusing on land value, aligning with the Ricardian theory that emphasizes the role of location and societal demand in generating land rents, thereby promoting equity and reducing speculation.
\item \textbf{Answer:} The relationship between the demand for stocks and bonds and the asset demand for money is inversely related. When the demand for stocks and bonds increases, investors are more likely to allocate their funds into these financial assets rather than holding cash, leading to a decrease in the asset demand for money. This shift occurs because individuals seek to maximize returns on their investments, favoring assets that are expected to appreciate in value over the liquidity provided by money. Consequently, as more capital flows into the stock and bond markets, the overall demand for money diminishes, reflecting a preference for earning potential over liquidity. Thus, fluctuations in the financial markets can significantly influence how much money individuals choose to hold as an asset.
\\ \textit{Note:} correct because it accurately describes the inverse relationship between the demand for stocks and bonds and the asset demand for money, illustrating how increased investment in financial assets diminishes the desire to hold cash due to the pursuit of higher returns.
\end{enumerate}

\vfill
\noindent\textit{End of Answer Key}
\end{document}